% =====================================================================
% Jacob E. Brinton — Curriculum Vitae  (GENERATED from resume_src/data.yml)
% DO NOT EDIT THIS FILE DIRECTLY. Edit resume_src/data.yml and rerun
% `python resume_src/build.py`. Hand-edits here will be clobbered.
% =====================================================================
\documentclass[11pt,letterpaper]{extarticle}

% ---------- Packages ----------
\usepackage[top=0.7in,bottom=0.7in,left=0.9in,right=0.9in]{geometry}
\usepackage[hidelinks]{hyperref}
\usepackage{enumitem}
\usepackage{titlesec}
\usepackage{xcolor}
\usepackage[T1]{fontenc}
\usepackage{lmodern}
\usepackage{microtype}
\usepackage{etoolbox}

% ---------- Hyperlink colors ----------
\hypersetup{
  colorlinks=true,
  urlcolor=black,
  linkcolor=black,
  citecolor=black,
}

\pagestyle{plain}
\setlength{\parindent}{0pt}
\setlength{\parskip}{2pt}

% ---------- Section formatting ----------
\titleformat{\section}
  {\normalfont\large\bfseries\MakeUppercase}
  {}{0em}{}
  [\vspace{2pt}\hrule height 0.4pt\relax]
\titlespacing*{\section}{0pt}{12pt}{6pt}

% ---------- List formatting ----------
\setlist[itemize]{
  leftmargin=1.4em,
  labelsep=0.5em,
  itemsep=1pt,
  parsep=0pt,
  topsep=2pt,
  partopsep=0pt,
  label=\small$\bullet$
}

% ---------- Entry macros ----------
% parbox-based layout: left column holds title/org (wraps if too long),
% right column holds date. Prevents title/date collisions when the title
% is longer than the text width (e.g. long paper titles in Publications).
\newcommand{\entry}[4]{%
  \par\vspace{3pt}%
  \noindent\begin{minipage}[t]{0.68\textwidth}\raggedright\textbf{#1}\end{minipage}%
  \hfill\begin{minipage}[t]{0.3\textwidth}\raggedleft\textit{#2}\end{minipage}\par
  \ifstrempty{#3}{%
    \ifstrempty{#4}{}{\noindent\hfill #4\par}%
  }{%
    \noindent\begin{minipage}[t]{0.68\textwidth}\raggedright\textit{#3}\end{minipage}%
    \ifstrempty{#4}{}{\hfill\begin{minipage}[t]{0.3\textwidth}\raggedleft #4\end{minipage}}\par
  }%
}

\newcommand{\projentry}[2]{%
  \par\vspace{3pt}%
  \noindent\begin{minipage}[t]{0.68\textwidth}\raggedright\textbf{#1}\end{minipage}%
  \hfill\begin{minipage}[t]{0.3\textwidth}\raggedleft\textit{#2}\end{minipage}\par
}

% ---------- CV header ----------
\newcommand{\cvheader}{%
  \begin{center}
    {\LARGE\bfseries Jacob E. Brinton}\\[5pt]
    \small
    \href{mailto:jbrin@bu.edu}{jbrin@bu.edu}
    \;$\cdot$\;
    \href{https://jebrinton.github.io}{jebrinton.github.io}
    \;$\cdot$\;
    Boston, Massachusetts
    \\[2pt]
    \href{https://linkedin.com/in/jbrinton}{linkedin.com/in/jbrinton}
    \;$\cdot$\;
    \href{https://github.com/jebrinton}{github.com/jebrinton}
  \end{center}
  \vspace{2pt}
}

% =====================================================================
\begin{document}
\cvheader

% ---------- EDUCATION ----------
\section{Education}
\entry{Boston University}{May 2026}{B.A./M.S. in Mathematics and Computer Science}{GPA: 3.89 (B.A.); 4.0 (M.S.)}
\begin{itemize}
  \item \textbf{Relevant coursework:} Interpretable Machine Learning, Directed Study in LLM Translation Interpretability, Principles of ML, Natural Language Processing, Stochastic Processes, Causal Inference, Abstract Algebra, Advanced Algorithms, Intro to AI, Foundations of Mathematics, Logic and Programming Languages, Theory of Computation\end{itemize}
\entry{Utah State University}{Aug 2022 -- May 2023}{}{GPA: 4.0}
\begin{itemize}
  \item \textbf{Relevant coursework:} Honors Calculus II \& III, Chemistry, Physics, Programming in C, CAD\end{itemize}

% ---------- RESEARCH & PROFESSIONAL EXPERIENCE ----------
\section{Research and Professional Experience}
\entry{Technical Fellowship Reviewer}{May 2026 -- Present}{Constellation}{Remote Contractor}
\begin{itemize}
  \item Screened and reference checked 200+ applicants to the Anthropic Fellows Program and the Astra Fellowship.
\end{itemize}
\entry{Technical Interviewer}{March 2026 -- Present}{FAR.AI}{Remote Contractor}
\begin{itemize}
  \item Interviewed research scientist and engineer candidates across a variety of ML, math, and programming questions.
\end{itemize}
\entry{Executive Director}{May 2025 -- Present}{Boston University AI Safety \& Alignment}{Boston, MA}
\begin{itemize}
  \item Increased membership by 10x; members accepted to ERA, LASR, MARS, SPAR; Harvard and MIT graduate programs.
  \item Presented on a wide range of AI risks to over a dozen courses, reaching a total enrollment of 1{,}000+ students.
  \item Negotiated hackathon partnerships; organized 20 events, from guest speakers to interuni. socials and research meetings.
  \item Wrote and facilitated technical AI safety curriculum; organized 6 discussion cohorts and 2 upskilling programs.
\end{itemize}
\entry{AI Safety Fieldbuilding Mentor}{Dec 2025 -- May 2026}{Pathfinder Fellowship (Kairos)}{Remote}
\begin{itemize}
  \item Guided Columbia University's AI safety group directors on outreach strategy, lecture series, and leadership transitions.
\end{itemize}
\entry{Research Fellow}{Dec 2025 -- Present}{Mentorship for Alignment Researchers (Cambridge AI Safety Hub)}{Cambridge, UK}
\begin{itemize}
  \item Developed model diffing interpretability methods for model safeguard checkpoints; mentored by Prof. Zhijing Jin.
\end{itemize}
\entry{Teaching Fellow, Course Assistant}{Aug 2024 -- Dec 2025}{Boston University College of Arts and Sciences, College of Engineering}{Boston, MA}
\begin{itemize}
  \item Taught in lab sections and office hours for linear algebra, computer systems, and calculus courses.
\end{itemize}
\entry{Machine Learning Research Intern}{May 2023 -- Aug 2023}{NSF Research Experiences for Undergraduates, Montana State University}{Bozeman, MT}
\begin{itemize}
  \item Designed UAV flight experiment, collecting 3 hours of audio; applied signal processing and ML methods for detection.
\end{itemize}
\entry{Undergraduate Researcher}{Dec 2022 -- May 2023}{DIRECT Laboratory (Utah State University)}{Logan, UT}
\begin{itemize}
  \item Wrote Python-based publishers and subscribers in ROS for robot pathfinding and camera movement.
\end{itemize}
\entry{Ground Station Subteam Member}{Sep 2022 -- May 2023}{Get Away Special Satellite Team}{Logan, UT}
\begin{itemize}
  \item 3D-modeled and manufactured prototypes in SolidWorks for an S-band patch antenna connection and coverglass frame.
\end{itemize}

% ---------- PUBLICATIONS ----------
\section{Publications}
\projentry{Assessing LWIR and Visual Sensor Modalities for UAV Detection and Tracking}{May 2024}
\begin{itemize}
  \item Coauthored paper assessing long-wave infrared (LWIR) and visual sensor modalities for UAV detection.
  \item \textit{2024 IEEE Intermountain Engineering, Technology and Computing (IETC), pp.~141--146}.\\ \href{https://doi.org/10.1109/IETC61393.2024.10564431}{doi:10.1109/IETC61393.2024.10564431}\end{itemize}

% ---------- PROJECTS (unpublished / ongoing / coursework) ----------
\section{Projects}
\projentry{The Interlingua Hypothesis: LLMs Translate via a Latent Task-agnostic Feature Space}{Mar 2025 -- Present}
\begin{itemize}
  \item Experimented on whether LLMs translate via shared multilingual representations as opposed to pairwise mappings.
  \item Used gradient-based methods to identify and then causally validate components implicated in translation.
  \item Quantified multilinguality of features' receptive fields and output spaces, and representation similarity across layers.
  \item Presented posters at the New England Mechanistic Interpretability Workshop and Boston University's UROP Symposium.
\end{itemize}
\projentry{Cross-Language Synonym Analysis and Scoring}{Sep 2024 -- Jan 2025}
\begin{itemize}
  \item Devised and empirically computed metrics for cross-language synonym closeness and translation flexibility.
\end{itemize}
\projentry{Tetris Reinforcement Learning Agent}{Apr 2024 -- May 2024}
\begin{itemize}
  \item Designed Q-learning agent with heuristic-guided exploration and reward shaping; placed 2nd in 150-person course.
\end{itemize}

% ---------- PRESENTATIONS & POSTERS ----------
\section{Presentations and Posters}
\par\vspace{2pt}%
\noindent\textbf{Poster: Noisy Translation Channels in Large Language Models}\hfill\textit{2025}\par
\noindent\textit{New England Mechanistic Interpretability workshop (NEMI)}\par

% ---------- AWARDS & FELLOWSHIPS ----------
\section{Awards and Fellowships}
\begin{itemize}
  \item \textbf{Boston University Undergraduate Research Opportunities Program (UROP) \textemdash\ Summer Research Fellowship} \hfill \textit{Summer 2025}\\ \small Funded the Noisy Translation Channels project; presented poster at NEMI.  \item \textbf{Tetris Reinforcement Learning Competition \textemdash\ 2nd of 150} \hfill \textit{May 2024}\\ \small CS 506 Tools for Data Science, Boston University.\end{itemize}

% ---------- SKILLS ----------
\section{Technical Skills and Languages}
\noindent\textbf{Technical:} Python, PyTorch, HuggingFace, nnsight, Weights \& Biases, Claude Code, OCaml, C.\par
\noindent\textbf{Languages:} English (Native), Spanish (Fluent).

\end{document}
